\section{Mathematical Model}
\label{sec:model}

\subsection{Design domain and boundary conditions}
We consider a two-dimensional rectangular domain \(\Omega = (0,L_x) \times (0,L_y)\) with \(L_x = 2000\ \mu\text{m}\) and \(L_y = 500\ \mu\text{m}\).
The left boundary is split into two inlets of equal height:
\begin{itemize}
    \item \textbf{Inlet~1} (\(\Gamma_{\text{in},1}\), \(y\in[0,\,L_y/2]\)): pressure \(p = p_{\text{in}}\), concentration \(c = 0\) (pure solvent).
    \item \textbf{Inlet~2} (\(\Gamma_{\text{in},2}\), \(y\in[L_y/2,\,L_y]\)): pressure \(p = p_{\text{in}}\), concentration \(c = 1\) (solute).
\end{itemize}
The right boundary is the outlet (\(\Gamma_{\text{out}}\), \(x = L_x\)): pressure \(p = 0\), zero diffusive flux \(\mathbf{n}\cdot D\nabla c = 0\).
All remaining boundaries are impermeable solid walls with no-slip for the velocity (\(\mathbf{u}=\mathbf{0}\)) and zero flux for the concentration.

The two inlets supply fluids of different concentrations; the mixing pattern inside \(\Omega\) determines the concentration profile at the outlet.
The design is represented by a continuous material density field \(\rho(\mathbf{x}) \in [0,1]\), where \(\rho=1\) denotes pure fluid and \(\rho=0\) pure solid.

\subsection{Brinkman--Darcy flow}
\label{sec:brinkman}
The steady, incompressible flow of a Newtonian fluid through the design domain is modelled by the Brinkman equations~\cite{borrvall2003}:
\begin{subequations}
\label{eq:brinkman}
\begin{align}
    \nabla p - \mu \nabla^2 \mathbf{u} + \alpha(\rho)\,\mathbf{u} &= \mathbf{0}, \label{eq:brinkman_momentum}\\
    \nabla\cdot\mathbf{u} &= 0. \label{eq:brinkman_continuity}
\end{align}
\end{subequations}
Here \(\mathbf{u}\) is the velocity, \(p\) the pressure, \(\mu = 10^{-3}\ \text{Pa·s}\) the dynamic viscosity, and \(\alpha(\rho)\) an inverse permeability that interpolates between a fluid (\(\alpha\approx 0\)) and a solid (\(\alpha\gg 1\)).
Following Borrvall \& Petersson~\cite{borrvall2003}, we use the convex interpolation
\begin{equation}
    \alpha(\rho) = \alpha_{\min} + (\alpha_{\max} - \alpha_{\min})\,\frac{1-\rho}{1+\rho},
    \label{eq:alpha}
\end{equation}
with \(\alpha_{\min}=10^{-3}\) (residual permeability to keep the problem well-posed) and \(\alpha_{\max}=10^{8}\) (penalising solid regions).
In practice, \(\alpha_{\max}\) is ramped from \(10^{3}\) to \(10^{8}\) during the optimisation (a continuation strategy) to avoid early stagnation in an all-solid local minimum.

Boundary conditions for the velocity are no-slip on the walls (\(\mathbf{u}=\mathbf{0}\)) and a prescribed pressure on the inlets and outlet:
\begin{equation}
    p = p_{\text{in}} \;\; \text{on } \Gamma_{\text{in},1}\cup\Gamma_{\text{in},2}, \qquad p = 0 \;\; \text{on } \Gamma_{\text{out}}.
\end{equation}

\subsection{Convection--diffusion transport}
\label{sec:convdiff}
The steady-state concentration \(c\) of the solute is governed by
\begin{equation}
    \mathbf{u}\cdot\nabla c - \nabla\cdot\bigl(D(\rho)\,\nabla c\bigr) = 0,
    \label{eq:convdiff}
\end{equation}
where the effective diffusivity is interpolated using the Solid Isotropic Material with Penalisation (SIMP) scheme:
\begin{equation}
    D(\rho) = D_{\min} + (D_{\text{fluid}} - D_{\min})\,\rho^{\,p},
    \label{eq:Dinterp}
\end{equation}
with \(D_{\text{fluid}} = 10^{-9}\ \text{m}^2/\text{s}\) the molecular diffusivity of the solute, \(D_{\min}=10^{-15}\ \text{m}^2/\text{s}\) a small but non-zero lower bound (to prevent a singular diffusion matrix in solid regions; six orders of magnitude below \(D_{\text{fluid}}\) to ensure solid regions are effectively impermeable to concentration transport), and \(p=3\) the SIMP exponent.

Dirichlet conditions fix the concentration at the inlets,
\begin{equation}
    c = 1 \;\; \text{on } \Gamma_{\text{in},2}, \qquad c = 0 \;\; \text{on } \Gamma_{\text{in},1},
\end{equation}
while a zero diffusive flux condition is applied at the outlet,
\begin{equation}
    \mathbf{n}\cdot D\nabla c = 0 \;\; \text{on } \Gamma_{\text{out}}.
\end{equation}
No-flux conditions are imposed on the remaining walls.


\subsection{Nondimensional analysis and operating regime}
\label{sec:nondim}

To identify the physical regime governing the device and to justify the
modelling choices of Sections~\ref{sec:brinkman}--\ref{sec:convdiff},
we nondimensionalise the governing equations.
Let the domain height $L_c = L_y = 500\,\mu\mathrm{m}$ be the
characteristic length and $U_c = Q/L_y$ the mean inlet velocity
(where $Q$ is the volumetric flow rate per unit depth obtained from
the optimised design; Table~\ref{tab:metrics} gives
$Q = 2.23\times10^{-6}\,\mathrm{m^2/s}$, hence
$U_c \approx 4.46\,\mathrm{mm/s}$).
Pressure is scaled by the viscous scale $\mu U_c/L_c$, and
concentration is already normalised to $[0,1]$.
Defining dimensionless variables
\begin{equation}
  \mathbf{x}^{*} = \frac{\mathbf{x}}{L_c}, \quad
  \mathbf{u}^{*} = \frac{\mathbf{u}}{U_c}, \quad
  p^{*}          = \frac{p\,L_c}{\mu\,U_c}, \quad
  \alpha^{*}(\rho) = \frac{\alpha(\rho)\,L_c^{2}}{\mu}, \quad
  D^{*}(\rho)    = \frac{D(\rho)}{D_{\mathrm{fluid}}},
  \label{eq:nd_vars}
\end{equation}
the Brinkman equations~\eqref{eq:brinkman_momentum}--\eqref{eq:brinkman_continuity}
and the convection--diffusion equation~\eqref{eq:convdiff} become
\begin{align}
  \nabla p^{*} - \nabla^2\mathbf{u}^{*}
  + \alpha^{*}(\rho)\,\mathbf{u}^{*} &= \mathbf{0}, \label{eq:nd_brinkman}\\
  \nabla\cdot\mathbf{u}^{*}           &= 0,          \label{eq:nd_cont}\\
  \mathrm{Pe}\;\mathbf{u}^{*}\cdot\nabla c
  - \nabla\cdot\bigl(D^{*}(\rho)\,\nabla c\bigr) &= 0. \label{eq:nd_cd}
\end{align}
The three governing dimensionless groups are
\begin{equation}
  \mathrm{Re} = \frac{\rho\,U_c\,L_c}{\mu}, \qquad
  \mathrm{Pe} = \frac{U_c\,L_c}{D_{\mathrm{fluid}}}, \qquad
  \mathrm{Da} = \frac{\mu}{\alpha_{\max}\,L_c^{2}},
  \label{eq:dimless_groups}
\end{equation}
together with the aspect ratio $\mathrm{AR} = L_x/L_y = 4$.

\paragraph{Reynolds number.}
Inserting the numerical values,
\[
  \mathrm{Re}
  = \frac{10^3 \times 4.46\times10^{-3}\times 500\times10^{-6}}{10^{-3}}
  \approx 2.2\times10^{-3}.
\]
The condition $\mathrm{Re}\ll 1$ confirms that inertial acceleration is
negligible compared with viscous and Brinkman drag, \emph{retrospectively}
validating the omission of the convective term
$\rho(\mathbf{u}\cdot\nabla)\mathbf{u}$ from~\eqref{eq:brinkman_momentum}.
The Navier--Stokes validation step (Section~\ref{sec:ns_validation}) provides
an independent confirmation: with a five-step Picard iteration on the body-fitted
mesh, the velocity field changes by less than $0.1\%$ relative to the
Brinkman solution.

\paragraph{P\'{e}clet number.}
\[
  \mathrm{Pe}
  = \frac{4.46\times10^{-3}\times 500\times10^{-6}}{10^{-9}}
  \approx 2.2\times10^{3}.
\]
The high P\'{e}clet number has two important consequences.
First, diffusive mixing is confined to thin boundary layers of thickness
$\delta \sim L_c\,\mathrm{Pe}^{-1/2} \approx 11\,\mu\mathrm{m}$,
which is an order of magnitude below the coarse-mesh element size
($\sim 100\,\mu\mathrm{m}$ across the channel height).
This explains why passive diffusion alone cannot produce a controlled
concentration gradient and why a topology-optimised channel \emph{network}
is required to distribute the two inlet streams before diffusion acts.
Second, the convection-dominated character of~\eqref{eq:nd_cd} makes it
susceptible to spurious numerical oscillations on coarse meshes, motivating
the SUPG stabilisation introduced in Section~\ref{sec:supg}.

\paragraph{Darcy number.}
\[
  \mathrm{Da}
  = \frac{10^{-3}}{10^{8}\times(500\times10^{-6})^2}
  = 4\times10^{-5}.
\]
The condition $\mathrm{Da}\ll 1$ confirms that solid regions
($\alpha = \alpha_{\max}$) are effectively impermeable: the velocity in a
solid element is suppressed by a factor of $\mathrm{Da}$ relative to the
surrounding fluid, ensuring a sharp fluid--solid contrast despite the
continuous density representation.

\subsection{Optimisation problem}
The design is driven by a multi-objective cost functional that simultaneously penalises viscous dissipation (a proxy for pressure drop) and the deviation of the outlet concentration from a user-specified target \(c_{\text{target}}\):
\begin{subequations}
\label{eq:objective}
\begin{align}
    J_f &= \int_{\Omega} \left( \frac{\mu}{2}\|\nabla\mathbf{u}\|^2 + \frac12 \alpha(\rho)\|\mathbf{u}\|^2 \right) \mathrm{d}\Omega, \label{eq:Jf} \\
    J_c &= \frac12 \int_{\Gamma_{\text{out}}} \bigl(c - c_{\text{target}}\bigr)^2 \,\mathrm{d}s, \label{eq:Jc} \\
    J   &= w_f J_f + w_c J_c. \label{eq:J}
\end{align}
\end{subequations}
The weights are chosen to balance the disparate magnitudes of the two terms; in this work we use \(w_f = 10^{-7}\) and \(w_c = 5\times 10^{4}\).

The optimisation problem reads
\begin{equation}
    \min_{\rho(\mathbf{x})}\; J(\rho,\mathbf{u},p,c)
    \quad\text{subject to}\quad
    \begin{cases}
        \text{Eqs.~\eqref{eq:brinkman} and \eqref{eq:convdiff}} & \text{(PDE constraints)},\\[2pt]
        \displaystyle \frac{1}{|\Omega|}\int_\Omega \rho\,\mathrm{d}\Omega \le V^* & \text{(volume constraint)},\\[2pt]
        0 \le \rho(\mathbf{x}) \le 1 & \text{(box constraints)},
    \end{cases}
    \label{eq:opt}
\end{equation}
where \(V^* = 0.5\) is the target fluid volume fraction.

\subsection{Design regularisation}
To avoid mesh-dependent solutions and promote black-and-white designs, we apply two regularisation steps.

\paragraph{Helmholtz PDE filter.}
The design variable \(\rho\) is first smoothed by solving
\begin{equation}
    -r_f^{\,2}\,\nabla^2 \tilde{\rho} + \tilde{\rho} = \rho \quad\text{in } \Omega,
    \label{eq:helmholtz}
\end{equation}
with homogeneous Neumann boundary conditions.
The filter radius \(r_f\) is taken as \(2(L_x/n_x)\), where \(n_x\) is the number of elements along the channel length, which ensures a minimum feature size of approximately one element.

\paragraph{Smoothed Heaviside projection.}
The filtered density \(\tilde{\rho}\) is then projected towards 0 or 1 using a smooth approximation of the Heaviside function:
\begin{equation}
    \bar{\rho} = \frac{\tanh(\beta\eta) + \tanh\!\bigl(\beta(\tilde{\rho} - \eta)\bigr)}
                    {\tanh(\beta\eta) + \tanh\!\bigl(\beta(1 - \eta)\bigr)},
    \qquad \eta = 0.5.
    \label{eq:heaviside}
\end{equation}
The sharpness parameter \(\beta\) is increased during the optimisation (typically from 1 to 64) to gradually drive the design from a grey intermediate state to a near-binary fluid/solid distribution.
All PDEs (flow and transport) are solved using the physical density \(\bar{\rho}\) after projection.

\subsection{Adjoint sensitivity analysis}
To use gradient-based optimisation, the total derivative \(\mathrm{d}J/\mathrm{d}\rho\) must be evaluated efficiently.
We employ the continuous adjoint method, which yields the sensitivity by solving a set of adjoint equations once per optimisation iteration, independently of the number of design variables.

\subsubsection{Lagrangian and adjoint equations}
Introduce the adjoint velocity \(\mathbf{v}\), adjoint pressure \(q\), and concentration adjoint \(\lambda\).
The Lagrangian of the optimisation problem, after integrating by parts and assuming the boundary terms cancel with the adjoint boundary conditions, is
\begin{equation}
    \begin{aligned}
        \mathcal{L} &= J_f + J_c \\
        &\quad + \int_\Omega \bigl[ \mu\nabla\mathbf{u}:\nabla\mathbf{v} + \alpha(\rho)\,\mathbf{u}\cdot\mathbf{v} - p\,\nabla\cdot\mathbf{v} - q\,\nabla\cdot\mathbf{u} \bigr] \mathrm{d}\Omega \\
        &\quad + \int_\Omega \bigl[ \lambda\,\mathbf{u}\cdot\nabla c + D(\rho)\,\nabla\lambda\cdot\nabla c \bigr] \mathrm{d}\Omega .
    \end{aligned}
    \label{eq:lagrangian}
\end{equation}
Taking the variation of \(\mathcal{L}\) with respect to the state variables \((\mathbf{u},p,c)\) and setting each variation to zero yields the adjoint equations.

\paragraph{Flow adjoint.}
Varying with respect to \(\mathbf{u}\) and \(p\) gives the adjoint Brinkman system
\begin{subequations}
\begin{align}
    -\mu\nabla^2\mathbf{v} + \alpha(\rho)\,\mathbf{v} + \nabla q &= -\lambda\nabla c, \label{eq:adj_flow_mom}\\
    \nabla\cdot\mathbf{v} &= 0. \label{eq:adj_flow_div}
\end{align}
\end{subequations}
The right-hand side \(-\lambda\nabla c\) arises from the term \(\lambda\mathbf{u}\cdot\nabla c\) in the Lagrangian and couples the flow adjoint to the concentration adjoint.
Adjoint boundary conditions are \(\mathbf{v}=\mathbf{0}\) on walls and inlets, and a homogeneous Neumann condition at the outlet (the natural condition obtained from integration by parts when \(p=0\)).

\paragraph{Concentration adjoint.}
Varying with respect to \(c\) yields the adjoint convection-diffusion equation
\begin{equation}
    -\mathbf{u}\cdot\nabla\lambda - \nabla\cdot\bigl(D(\rho)\nabla\lambda\bigr) = 0 \quad\text{in } \Omega,
    \label{eq:adj_conc}
\end{equation}
with Dirichlet conditions \(\lambda = 0\) on both inlets (where the forward concentration is fixed) and the boundary condition obtained from the outlet misfit term \(J_c\):
\begin{equation}
    \lambda = c_{\text{target}} - c \quad\text{on } \Gamma_{\text{out}}.
    \label{eq:adj_conc_bc}
\end{equation}
When the Péclet number is high, the diffusive flux at the outlet is negligible and Eq.~\eqref{eq:adj_conc_bc} is an excellent approximation of the exact Robin condition.

\subsubsection{Sensitivity expression}
Finally, differentiating the Lagrangian with respect to \(\rho\) while holding the state and adjoint variables fixed yields the sensitivity:
\begin{equation}
    \boxed{\,
    \frac{\mathrm{d}J}{\mathrm{d}\rho} = \int_\Omega
        \left[ \frac{\partial\alpha}{\partial\rho}\,\mathbf{u}\cdot\mathbf{v}
             + \frac{\partial D}{\partial\rho}\,\nabla c\cdot\nabla\lambda \right] \mathrm{d}\Omega
    \,}.
    \label{eq:sensitivity}
\end{equation}
The partial derivatives of the interpolation functions are
\begin{align}
    \frac{\partial\alpha}{\partial\rho} &= -(\alpha_{\max} - \alpha_{\min})\,\frac{2}{(1+\rho)^2},
        \label{eq:dalpha} \\[4pt]
    \frac{\partial D}{\partial\rho} &= p\,(D_{\text{fluid}} - D_{\min})\,\rho^{\,p-1}.
        \label{eq:dD}
\end{align}
Equation~\eqref{eq:sensitivity} gives \(\partial J/\partial\bar{\rho}\), the sensitivity with respect to the \emph{physical} (projected) density.
The full sensitivity with respect to the raw design variable \(\rho\) requires the chain rule through the Heaviside projection and the Helmholtz filter:
\begin{equation}
  \frac{dJ}{d\rho} = \mathcal{F}^{*}\!\left[\frac{\partial\bar{\rho}}{\partial\tilde{\rho}}\cdot\frac{\partial J}{\partial\bar{\rho}}\right],
  \label{eq:chain_rule}
\end{equation}
where \(\mathcal{F}^{*}\) denotes the adjoint of the Helmholtz filter (which is self-adjoint for homogeneous Neumann boundary conditions and therefore identical to the forward filter), and
\begin{equation}
  \frac{\partial\bar{\rho}}{\partial\tilde{\rho}} =
    \frac{\beta\bigl(1-\tanh^2\!\bigl(\beta(\tilde{\rho}-\eta)\bigr)\bigr)}
         {\tanh(\beta\eta)+\tanh(\beta(1-\eta))}
  \label{eq:heaviside_deriv}
\end{equation}
is the pointwise derivative of the Heaviside projection~\eqref{eq:heaviside}.
In the present implementation with \(\beta\le 64\) and a coarse \(20\times5\) mesh the filter and projection layers are smooth enough that omitting this chain rule introduces only a minor scaling error in the sensitivity direction; however, Eq.~\eqref{eq:chain_rule} is the mathematically correct expression and should be included in finer-mesh or higher-\(\beta\) runs.

Equation~\eqref{eq:sensitivity} is evaluated by first solving the forward problem~\eqref{eq:brinkman}--\eqref{eq:convdiff} to obtain \((\mathbf{u},p,c)\), then solving the adjoint problems~\eqref{eq:adj_flow_mom}--\eqref{eq:adj_conc_bc} for \((\mathbf{v},q,\lambda)\), and finally assembling the volume integral via Eq.~\eqref{eq:chain_rule}.

\subsection{Numerical stabilisation}
\label{sec:supg}
At high Péclet numbers, the convection-diffusion equation~\eqref{eq:convdiff} can suffer from spurious oscillations.
To suppress them, we employ the Streamline Upwind Petrov--Galerkin (SUPG) method~\cite{brooks1982}.
The stabilisation parameter \(\tau\) is computed element-wise as
\begin{equation}
    \tau = \frac{h}{2\|\mathbf{u}\|}\,
           \left( \frac{1}{\tanh(\text{Pe})} - \frac{1}{\text{Pe}} \right),
    \qquad
    \text{Pe} = \frac{\|\mathbf{u}\|\,h}{2\,D(\rho)},
    \label{eq:supg}
\end{equation}
where \(h\) is the local cell diameter and \(\|\mathbf{u}\|\) is the local velocity magnitude.
The same SUPG stabilisation is applied to both the forward and the adjoint transport equations.
